\documentclass{article}
\usepackage[utf8]{inputenc}
\usepackage{amsmath}
\usepackage{geometry}
\geometry{margin=2cm}
\begin{document}

\section*{Vacinas e Resposta Imune – Questão 99 ENEM 2024 – Ciências da Natureza}

\subsection*{Interpretação e Estratégia}

O enunciado explica a importância da proteína espícula do coronavírus como antígeno nas vacinas contra a covid-19. A questão pede para identificar qual é o efeito dessa proteína viral no organismo vacinado.

\textbf{Termos-chave:}
\begin{itemize}
  \item \textbf{Induzir:} Provocar ou estimular uma resposta biológica, neste caso a resposta imunológica ativa.
  \item \textbf{Proteína espícula:} Estrutura antigênica que desencadeia a resposta imune nas vacinas contra o coronavírus.
\end{itemize}

\textbf{O que a questão pede:} Compreender o papel da proteína viral apresentada nas vacinas — a proteína espícula — e o efeito imunológico que ela induz após a vacinação.

\textbf{Estratégia:} Relacionar o conceito de antígeno ao estímulo para produção de anticorpos, reconhecer que vacinas induzem imunidade ativa e eliminar alternativas incorretas que descrevem ações erradas como imunidade passiva, alteração do DNA ou ação farmacológica direta da vacina.

\subsection*{Conteúdo e Mini Revisão}

\begin{center}
\begin{tabular}{|l|p{9cm}|}
\hline
\textbf{Conceito} & \textbf{Chave para a resposta} \\ \hline
\textbf{Antígeno} & Substância que induz resposta imunológica. A proteína espícula atua como antígeno na vacina, estimulando produção de anticorpos específicos. \\ \hline
\textbf{Imunidade ativa} & Produção interna de anticorpos pelo organismo. Vacinas estimulam a produção própria de anticorpos específicos. \\ \hline
\textbf{Imunidade passiva} & Recebimento direto de anticorpos prontos, não ocorre neste caso. \\ \hline
\textbf{Neutralização viral} & Bloqueio da ação do vírus por anticorpos. Resultado da resposta imunológica, não da vacina diretamente. \\ \hline
\end{tabular}
\end{center}

\textbf{Teoria:}
\begin{itemize}
  \item As vacinas usam antígenos para induzir resposta imune específica. No caso da covid-19, a proteína espícula é esse antígeno.
  \item Imunidade vacinal é ativa, pois induz o organismo a produzir anticorpos e células de memória.
  \item Vacinas não alteram o material genético ou os receptores celulares.
\end{itemize}

\subsection*{Resolução e Pulo do Gato}

Sabemos que a vacina apresenta a proteína espícula como antígeno, promovendo a produção ativa de anticorpos específicos.

Assim, a alternativa \textbf{A} está correta: a vacina induz a produção de anticorpos contra a proteína viral, preparando o sistema imunológico para combater o vírus.

As demais alternativas apresentam conceitos errados como imunidade passiva (B), alteração do genoma (C), ação direta da vacina na neutralização (D) ou modificação dos receptores celulares (E).

\textbf{Pulo do gato:} Lembre-se de que vacinas estimulam o corpo a produzir seus próprios anticorpos – imunidade ativa – e não fornecem anticorpos prontos nem alteram o DNA.

\subsection*{Armadilhas}

\begin{itemize}
  \item Confundir imunidade ativa com passiva.
  \item Achar que a vacina altera o DNA para obter proteção.
  \item Pensar que a vacina age diretamente no vírus, sem mediação imunológica.
  \item Imaginar a modificação dos receptores celulares pela vacina.
  \item Interpretação errada de ação imediata da vacina, em vez de indução da resposta imune.
\end{itemize}

\subsection*{Código da Questão}

\texttt{CN\_BIOLOGIA\_MOLECULAR\_001}

\clearpage


\end{document}